\documentclass[a4paper,10pt]{article}

\usepackage[T1]{fontenc}
\usepackage[utf8]{inputenc}
\usepackage[brazil]{babel}
\usepackage{lmodern}
\usepackage[margin=1.55cm]{geometry}
\usepackage{enumitem}
\usepackage{tabularx}
\usepackage{xcolor}
\usepackage[hidelinks]{hyperref}
\usepackage{titlesec}
\usepackage{microtype}
\usepackage{parskip}

\definecolor{accent}{HTML}{1F4E79}
\hypersetup{colorlinks=true, urlcolor=accent}
\pagestyle{empty}
\setlength{\parindent}{0pt}
\setlength{\parskip}{2.5pt}
\setlist[itemize]{leftmargin=1.25em, topsep=2pt, itemsep=1.5pt, parsep=0pt}
\titleformat{\section}{\large\bfseries\scshape\color{accent}}{}{0em}{}[\vspace{-3pt}\titlerule]
\titlespacing*{\section}{0pt}{9pt}{4pt}
\renewcommand{\arraystretch}{1.05}

\newcommand{\cvrole}[4]{%
  \begin{tabularx}{\linewidth}{@{}X r@{}}
    \textbf{#1} \enspace|\enspace \textit{#2} & \textbf{#3} \\
    \multicolumn{2}{@{}l@{}}{\textit{#4}}
  \end{tabularx}\vspace{-5pt}
}
\newcommand{\cvedu}[3]{%
  \begin{tabularx}{\linewidth}{@{}X r@{}}
    \textbf{#1} & \textbf{#2} \\
    \multicolumn{2}{@{}l@{}}{#3}
  \end{tabularx}\vspace{-4pt}
}

\begin{document}

\begin{center}
  {\LARGE\bfseries PAULO HENRIQUE}\\[3pt]
  {\large Engenheiro de Software}\\[5pt]
  \href{mailto:pauloh1288@gmail.com}{pauloh1288@gmail.com}
  \quad|\quad (85) 99951-1965
  \quad|\quad \href{https://github.com/youphenrique}{github.com/youphenrique}\\[-1pt]
  \href{https://www.linkedin.com/in/phenrique7}{linkedin.com/in/phenrique7}
  \quad|\quad \href{https://www.phenrique.me}{phenrique.me}
\end{center}

\section{Resumo profissional}
Engenheiro de Software com mais de 9 anos de experiência em JavaScript e TypeScript, especializado em aplicações web escaláveis, modernização de sistemas legados, integrações com APIs e automação de entrega. Atuação em produtos de alto tráfego, com mais de 10 mil visitas diárias e aplicações utilizadas por mais de 100 mil usuários. Experiência em arquitetura frontend, desenvolvimento full stack, qualidade de software, CI/CD e mentoria técnica.

\section{Competências técnicas}
\textbf{Linguagens:} JavaScript, TypeScript, Kotlin, PHP.\\
\textbf{Frontend:} React, Next.js, Vue.js, Chakra UI, Storybook, HTML, CSS, i18next, gerenciamento de estado, extensões de navegador.\\
\textbf{Backend e dados:} Node.js, Spring Boot, Laravel, Hapi.js, Express.js, GraphQL, REST APIs, Postgres, MySQL, MongoDB, Redis, BullMQ, Strapi.\\
\textbf{Cloud, DevOps e qualidade:} Docker, Docker Compose, GitHub Actions, Jenkins, SonarQube, AWS, Oracle Cloud, Coolify, Vitest, Playwright, Cypress, Jest, React Testing Library.

\section{Experiência profissional}
\cvrole{Nordeste Painéis}{Engenheiro de Software}{2025 -- Presente}{Híbrido}
\begin{itemize}
  \item Responsável pelo design, arquitetura e desenvolvimento de ERP customizado para gerenciamento de montagem e venda de quadros elétricos.
  \item Desenvolve interfaces responsivas com React e Chakra UI e APIs escaláveis com Spring Boot (Kotlin), além de automações em Node.js.
  \item Gerencia ambientes com Docker, pipelines de CI/CD no GitHub Actions e serviços em AWS e Oracle Cloud, utilizando Coolify para deploy.
\end{itemize}

\cvrole{T10}{Engenheiro de Software Frontend}{2020 -- 2024}{Remoto}
\begin{itemize}
  \item Desenvolveu aplicações e sites com TypeScript, React, Next.js, GraphQL, Storybook, Vitest e Playwright; criou hooks, fluxos de internacionalização e componentes modulares.
  \item Contribuiu para serviços backend em Node.js, Postgres, Strapi, Redis, BullMQ e Docker, fortalecendo segurança, escalabilidade e confiabilidade da stack.
  \item Otimizou produtos de alto tráfego com mais de 20 mil visitas diárias e aplicações usadas por mais de 100 mil usuários, mantendo desempenho e confiabilidade sob carga.
  \item Automatizou pipelines de CI/CD com GitHub Actions, Jenkins e SonarQube e mentorou profissionais juniores e plenos por meio de revisões de código e orientação técnica.
\end{itemize}

\cvrole{Secrelnet}{Engenheiro de Software}{2019 -- 2020}{Presencial}
\begin{itemize}
  \item Participou da migração de plataforma de backup em nuvem, de monólito PHP/Yii e MySQL para cliente React/TypeScript e API REST em Laravel, Redis e MySQL.
  \item Automatizou coleta de dados com Puppeteer e validou jornadas críticas com Cypress, Jest e React Testing Library; estruturou ambientes com Docker Compose.
  \item Desenvolveu sistema de gestão de monitores com React, Node.js, Hapi.js, Socket.IO e MongoDB, incluindo extensão de navegador em Vue.js e Webpack.
\end{itemize}

\newpage
\cvrole{Secrelnet}{Estagiário de Desenvolvimento de Software}{2017 -- 2019}{Presencial}
\begin{itemize}
  \item Manteve sistema PBX IP baseado em Asterisk, Yii, jQuery e MySQL; integrou e otimizou APIs Node.js, ambientes Docker e rotinas de servidor em Bash.
  \item Desenvolveu chat interno em React, Redux e Socket.IO integrado a API Node.js/Express, MySQL e comunicação em tempo real por WebSocket.
  \item Contribuiu para a migração de sistema de gestão de tarefas entre versões do Yii, com jQuery no cliente e SQL Server no armazenamento.
\end{itemize}

\cvrole{Universidade Estadual do Ceará (UECE)}{Mentor de Programação Acadêmica}{2014 -- 2017}{Presencial}
\begin{itemize}
  \item Realizou atividades de pesquisa e ensino para preparação em competições da Sociedade Brasileira de Computação (SBC), criando materiais de estudo e mentorando estudantes.
  \item Competiu na Maratona de Programação da SBC em 2014 e 2015; foi campeão estadual em 2016 e alcançou o 25\textsuperscript{o} lugar nacional entre 58 instituições.
\end{itemize}

\section{Educação}
\cvedu{Faculdade Cidade Teológica Pentecostal}{2023 -- 2025}{Pós-graduação em Sagradas Escrituras}
\cvedu{Universidade Estadual do Ceará}{2012 -- 2019}{Bacharelado em Ciência da Computação}
\small \textit{Fundamentos:} estruturas de dados e algoritmos, sistemas operacionais, redes, sistemas distribuídos, concorrência, padrões de projeto, arquitetura de sistemas, linguagens formais e compiladores.\normalsize

\section{Destaques adicionais}
\textbf{Áreas de atuação:} arquitetura de aplicações web, refatoração e migração de legados, performance, automação, testes, integração contínua, entrega contínua e mentoria técnica.\\
\textbf{Conquistas acadêmicas:} campeão estadual do Ceará na Maratona de Programação SBC (2016); finalista nacional pela UECE.

\end{document}
